\chapter{Khởi Động Bằng Câu Nói Vui}
\chapterintro{Vui vừa đủ để kéo lớp vào bài}{Một câu nói đúng độ sẽ làm lớp bật lên, nhưng vẫn phải giữ đường nối rõ ràng vào nội dung.}

\section{Thân vào lớp rồi, não vào chưa?}
\begin{khoidongbox}{Hoạt động khởi động}
Mở bằng câu hỏi dí dỏm: \textit{``Thân vào lớp rồi, não vào chưa?''}
\end{khoidongbox}
\begin{visaohieuqua}
Câu này đúng cảm giác của nhiều lớp sau giờ ra chơi. Học sinh thường bật cười và chấp nhận quay về tiết học dễ hơn.
\end{visaohieuqua}
\begin{cachlam5phut}
Nói xong thì nối luôn bằng một câu rất cụ thể như: \textit{``Rồi, kéo não vào bằng ý đầu tiên này.''}
\end{cachlam5phut}
\ngancach

\section{Hôm nay ai mở tiết giúp thầy cô?}
\begin{khoidongbox}{Hoạt động khởi động}
Mời lớp tìm một người mở tiết bằng một câu tóm, một ví dụ, hoặc một dự đoán vui.
\end{khoidongbox}
\begin{visaohieuqua}
Khi trao quyền mở tiết cho học sinh, lớp vào vai chủ động nhanh hơn. Không khí cũng bớt cảm giác bị kéo vào bài.
\end{visaohieuqua}
\begin{cachlam5phut}
Chỉ cần 1 em hoặc 1 bàn là đủ. Thầy cô chốt gọn rồi đi tiếp để giữ nhịp.
\end{cachlam5phut}
\ngancach

\section{Bài này nhìn hiền hay nhìn khó?}
\begin{khoidongbox}{Hoạt động khởi động}
Viết tên bài lên bảng rồi hỏi cả lớp: \textit{``Nhìn tên thôi, bài này hiền hay khó?''}
\end{khoidongbox}
\begin{visaohieuqua}
Học sinh thích những câu hỏi cảm tính đầu tiết vì dễ trả lời. Từ chỗ đó, thầy cô có thể dẫn vào mục tiêu bài rất mượt.
\end{visaohieuqua}
\begin{cachlam5phut}
Cho lớp giơ tay chọn, mời 2 em nói lý do, rồi nói luôn: \textit{``Mình xem nó hiền hay khó thật nhé.''}
\end{cachlam5phut}
\ngancach

\section{Nếu tiết học này có nhạc nền, đó là gì?}
\begin{khoidongbox}{Hoạt động khởi động}
Hỏi vui: \textit{``Nếu tiết học này có nhạc nền, em chọn loại nào?''}
\end{khoidongbox}
\begin{visaohieuqua}
Câu hỏi hình ảnh và gần đời sống tạo không khí nhẹ mà không cần đạo cụ. Nó rất hợp với lớp còn ngái ngủ hoặc ngại phát biểu.
\end{visaohieuqua}
\begin{cachlam5phut}
Lấy 2 đến 3 đáp án, rồi dùng một đáp án hay nhất để gọi tên tinh thần tiết học hôm nay.
\end{cachlam5phut}
\ngancach

\section{Một câu chào khác thường}
\begin{khoidongbox}{Hoạt động khởi động}
Thay lời chào quen bằng một câu lạ như: \textit{``Chào lớp học tử tế nhất trong 5 phút tới.''}
\end{khoidongbox}
\begin{visaohieuqua}
Lời chào khác thường làm lớp giật nhẹ sự chú ý. Điều quan trọng là câu đó vẫn chứa kỳ vọng tích cực và rõ ràng.
\end{visaohieuqua}
\begin{cachlam5phut}
Dùng 1 câu thôi, tránh diễn quá dài. Sau câu chào, chuyển thẳng vào yêu cầu đầu tiên của tiết học.
\end{cachlam5phut}